\section{Introdução}

A coleta eficiente de dados em redes de sensores sem fio (WSNs, do inglês \textit{wireless sensor networks}) é um componente essencial para a operação de sistemas de monitoramento ambiental, infraestrutura inteligente e aplicações em Internet das Coisas (IoT). Nessas redes, dispositivos de baixo consumo energético são distribuídos em uma área de monitoramento para coletar grandezas físicas (temperatura, umidade, vibração, entre outras) e transmiti-las periodicamente a um ponto central de coleta. Em cenários nos quais a conectividade com infraestrutura fixa é limitada ou inexistente, veículos aéreos não tripulados (UAVs, do inglês \textit{unmanned aerial vehicles}) têm sido explorados como coletores móveis, oferecendo flexibilidade de deslocamento, escalabilidade e maior cobertura espacial \cite{liu2018uavWSN}.

A utilização de UAVs como \emph{sinks} móveis — coletores de dados que absorvem as informações transmitidas pelos sensores da rede — introduz desafios complexos de planejamento de voo, coordenação e gerenciamento energético. O desempenho da coleta está diretamente relacionado a fatores como tempo de voo, energia disponível, capacidade de comunicação ar–terra e a estratégia de sobrevoo adotada. Três estratégias principais são identificadas na literatura \cite{messaoudi2023survey}: \textit{hovering} (o UAV paira estático sobre o sensor durante a coleta), \textit{fly-by} (o drone sobrevoa em movimento contínuo, coletando dados em trânsito) e híbrida (combinação das duas anteriores). A avaliação desse tipo de distema envolve, além de métricas clássicas (throughput agregado, latência, capacidade por nó e consumo energético), a \emph{Age of Information} (AoI), que mede o tempo desde a última atualização recebida de um sensor até o presente momento, crucial quando dados desatualizados perdem rapidamente relevância~\cite{messaoudi2023survey}.

A literatura recente propõe modelagens e técnicas para planejamento de trajetórias, incluindo (i) partição de área em células para coleta periódica \cite{liu2018uavWSN}, (ii) esquemas de coleta baseados em aprendizado por reforço profundo \cite{zhang2018drlrvc} e (iii) estratégias cooperativas entre múltiplos UAVs visando aumentar capacidade e eficiência energética \cite{yang2020surveyUAVIoT}. Em particular, o uso coordenado de múltiplos UAVs, quando bem modelado, pode proporcionar ganhos expressivos em capacidade por nó e no tempo de vida da rede em relação a abordagens com UAV único \cite{liu2018uavWSN,yang2020surveyUAVIoT}. Embora esses resultados motivem abordagens multiagente, este trabalho concentra-se no cenário com um único UAV, de modo a isolar o impacto das decisões de trajetória e consumo energético sobre as métricas de qualidade da coleta e viabilizar uma formulação de otimização exata tratável.

Neste trabalho, propomos uma formulação geral do problema de planejamento de voo para UAVs atuando como coletores de dados em WSNs, com foco em aspectos energéticos e operacionais. A modelagem visa abranger diferentes parâmetros envolvidos na coleta, como altura de voo, número de sensores, cobertura por célula e capacidade de coleta por unidade de tempo. O objetivo é otimizar simultaneamente métricas de eficiência energética, cobertura, tempo de coleta e AoI, oferecendo uma base para análise de desempenho e desenvolvimento de algoritmos de otimização aplicáveis a diferentes cenários de IoT assistida por UAVs.

O restante deste documento está organizado da seguinte forma. A Seção~2 apresenta os trabalhos relacionados e um comparativo com a proposta deste trabalho. A Seção~3 aborda a fundamentação teórica, incluindo conceitos de WSN, AoI, UAVs e modelos de energia e comunicação. A Seção~4 define os objetivos do trabalho. A Seção~5 descreve a modelagem matemática do problema de otimização. A Seção~6 detalha a metodologia experimental e os resultados obtidos para os dois modelos de drone avaliados. Por fim, a Seção~7 apresenta o cronograma de atividades.